
\documentclass{article}
\usepackage{colortbl}
\usepackage{makecell}
\usepackage{multirow}
\usepackage{supertabular}

\begin{document}

\newcounter{utterance}

\centering \large Interaction Transcript for game `hot\_air\_balloon', experiment `air\_balloon\_survival\_en\_negotiation\_hard', episode 3 with deepseek{-}chat{-}v3.1{-}t1.0.
\vspace{24pt}

{ \footnotesize  \setcounter{utterance}{1}
\setlength{\tabcolsep}{0pt}
\begin{supertabular}{c@{$\;$}|p{.15\linewidth}@{}p{.15\linewidth}p{.15\linewidth}p{.15\linewidth}p{.15\linewidth}p{.15\linewidth}}
    \# & $\;$A & \multicolumn{4}{c}{Game Master} & $\;\:$B\\
    \hline

    \theutterance \stepcounter{utterance}  
    & & \multicolumn{4}{p{0.6\linewidth}}{
        \cellcolor[rgb]{0.9,0.9,0.9}{
            \makecell[{{p{\linewidth}}}]{
                \texttt{\tiny{[P1$\langle$GM]}}
                \texttt{Stai partecipando a un gioco di negoziazione collaborativa.} \\
\\ 
\texttt{Insieme a un altro partecipante, dovete accordarvi su un unico set di oggetti da conservare. Ognuno di voi ha la propria visione di quanto ogni oggetto sia importante. Non sai come l'altro partecipante valuta gli oggetti. Inoltre, ti viene fornito lo sforzo che ogni oggetto richiede.} \\
\texttt{Potete accordarvi su un set solo se lo sforzo totale degli oggetti selezionati non supera un limite condiviso:} \\
\\ 
\texttt{LIMITE = 3681} \\
\\ 
\texttt{Ecco i valori di sforzo dei singoli oggetti:} \\
\\ 
\texttt{Sforzo oggetto = \{"C17": 700, "B41": 586, "A71": 505, "C00": 116, "B35": 664, "A77": 940, "A03": 387, "A04": 392, "C54": 209, "C01": 571, "A87": 4, "B32": 285, "B86": 651, "B33": 613, "A10": 740\}} \\
\\ 
\texttt{Ecco la tua visione personale sull'importanza di ogni oggetto:} \\
\\ 
\texttt{Valori di importanza dell'oggetto = \{"C17": 138, "B41": 583, "A71": 868, "C00": 822, "B35": 783, "A77": 65, "A03": 262, "A04": 121, "C54": 508, "C01": 780, "A87": 461, "B32": 484, "B86": 668, "B33": 389, "A10": 808\}} \\
\\ 
\texttt{Obiettivo:} \\
\\ 
\texttt{Il tuo obiettivo è negoziare un set condiviso di oggetti che ti avvantaggi il più possibile (cioè, massimizzi l'importanza totale per TE), rimanendo entro il LIMITEE. Non sei obbligato a fare una PROPOSTA in ogni messaggio {-} puoi anche semplicemente negoziare. Tutte le tattiche sono permesse!} \\
\\ 
\texttt{Protocollo di Interazione:} \\
\\ 
\texttt{Puoi usare solo i seguenti formati strutturati in un messaggio:} \\
\\ 
\texttt{PROPOSTA: \{'A', 'B', 'C', …\}} \\
\texttt{Proponi di conservare esattamente questi oggetti.} \\
\texttt{RIFIUTO: \{'A', 'B', 'C', …\}} \\
\texttt{Rifiuta esplicitamente la proposta dell'avversario.} \\
\texttt{ARGOMENTAZIONE: \{'...'\}} \\
\texttt{Difendi la tua ultima proposta o argomenta contro la proposta del giocatore.} \\
\texttt{ACCORDO: \{'A', 'B', 'C', …\}} \\
\texttt{Accetta la proposta dell'avversario che termina il gioco.} \\
\texttt{RAGIONAMENTO STRATEGICO: \{'...'\}} \\
\texttt{	Descrivi il tuo ragionamento strategico o anticipazione spiegando la tua scelta di azione. Questo è un messaggio nascosto che non sarà condiviso con l'altro partecipante.} \\
\\ 
\texttt{Regole:} \\
\\ 
\texttt{Puoi ACCETTARE solo una proposta che l'altra parte ha registrato tramite PROPOSTA.} \\
\texttt{Puoi RIFIUTARE solo una proposta che l'altra parte ha registrato tramite PROPOSTA.} \\
\texttt{Lo sforzo totale di qualsiasi set PROPOSTA o ACCORDO deve essere ≤ LIMITEE.} \\
\texttt{NON rivelare i tuoi punteggi di importanza nascosti.} \\
\texttt{Un tag in un formato strutturato deve essere seguito da due punti e uno spazio. L'argomento deve essere un set python contenente 0 o più stringhe.} \\
\texttt{Quindi, deve essere della forma TAG: \{...\}} \\
\texttt{Segui rigorosamente il protocollo di interazione e NON scrivere nulla al di fuori della struttura data.} \\
\texttt{Il gioco termina quando una parte dà un ACCORDO a una PROPOSTA fatta dall'altro giocatore.} \\
\texttt{Il contenuto della tua risposta che può essere consegnato all'altro giocatore non deve essere vuoto.} \\
\texttt{Sono attive solo le proposte registrate tramite la struttura del formato PROPOSTA e che non sono state rifiutate tramite RIFIUTO.} \\
\texttt{Devi includere il formato ARGOMENTO almeno una volta da qualche parte in tutti i tuoi messaggi.} \\
\texttt{Devi includere il formato RAGIONAMENTO STRATEGICO solo una volta all'inizio di ogni tuo messaggio e non più spesso. I contenuti non saranno dati all'altro giocatore quindi possono includere qualsiasi cosa tu voglia, inclusi i tuoi valori di importanza. Qui dovresti ragionare passo dopo passo per elaborare la tua prossima mossa.} \\
            }
        }
    }
    & \\ \\

    \theutterance \stepcounter{utterance}  
    & \multicolumn{4}{p{0.6\linewidth}}{
        \cellcolor[rgb]{0.9,0.9,0.9}{
            \makecell[{{p{\linewidth}}}]{
                \texttt{\tiny{[P1$\rangle$GM]}}
                \texttt{RAGIONAMENTO STRATEGICO: \{'Devo massimizzare il mio valore totale (8087) rimanendo sotto il limite di 3681. Analizzo gli oggetti con il miglior rapporto valore/sforzo per me: A87 (115.25), C00 (7.09), A71 (1.72), C54 (2.43), B32 (1.70), C01 (1.37), B86 (1.03), B35 (1.18), A10 (1.09), B41 (0.99). Gli oggetti con rapporto più basso sono A77 (0.07) e A04 (0.31). L\textbackslash{}'altro giocatore potrebbe avere preferenze diverse, quindi dovrei iniziare con una proposta forte che includa i miei oggetti ad alto valore/basso sforzo.'\}} \\
\\ 
\texttt{PROPOSTA: \{'A87', 'C00', 'A71', 'C54', 'B32', 'C01', 'B86', 'B35', 'A10', 'B41'\}} \\
\texttt{ARGOMENTAZIONE: \{'Questa proposta include oggetti ad alto valore con sforzo efficiente, massimizzando il beneficio complessivo mentre rimaniamo ben sotto il limite di sforzo. È una soluzione win{-}win che dovrebbe soddisfare entrambe le parti.'\}} \\
            }
        }
    }
    & & \\ \\

    \theutterance \stepcounter{utterance}  
    & & \multicolumn{4}{p{0.6\linewidth}}{
        \cellcolor[rgb]{0.9,0.9,0.9}{
            \makecell[{{p{\linewidth}}}]{
                \texttt{\tiny{[GM$\rangle$P2]}}
                \texttt{Stai partecipando a un gioco di negoziazione collaborativa.} \\
\\ 
\texttt{Insieme a un altro partecipante, dovete accordarvi su un unico set di oggetti da conservare. Ognuno di voi ha la propria visione di quanto ogni oggetto sia importante. Non sai come l'altro partecipante valuta gli oggetti. Inoltre, ti viene fornito lo sforzo che ogni oggetto richiede.} \\
\texttt{Potete accordarvi su un set solo se lo sforzo totale degli oggetti selezionati non supera un limite condiviso:} \\
\\ 
\texttt{LIMITE = 3681} \\
\\ 
\texttt{Ecco i valori di sforzo dei singoli oggetti:} \\
\\ 
\texttt{Sforzo oggetto = \{"C17": 700, "B41": 586, "A71": 505, "C00": 116, "B35": 664, "A77": 940, "A03": 387, "A04": 392, "C54": 209, "C01": 571, "A87": 4, "B32": 285, "B86": 651, "B33": 613, "A10": 740\}} \\
\\ 
\texttt{Ecco la tua visione personale sull'importanza di ogni oggetto:} \\
\\ 
\texttt{Valori di importanza dell'oggetto = \{"A71": 65, "C00": 121, "A10": 138, "B35": 262, "C01": 389, "B86": 461, "B41": 484, "C54": 508, "B32": 583, "A87": 668, "B33": 780, "A03": 783, "C17": 808, "A04": 822, "A77": 868\}} \\
\\ 
\texttt{Obiettivo:} \\
\\ 
\texttt{Il tuo obiettivo è negoziare un set condiviso di oggetti che ti avvantaggi il più possibile (cioè, massimizzi l'importanza totale per TE), rimanendo entro il LIMITEE. Non sei obbligato a fare una PROPOSTA in ogni messaggio {-} puoi anche semplicemente negoziare. Tutte le tattiche sono permesse!} \\
\\ 
\texttt{Protocollo di Interazione:} \\
\\ 
\texttt{Puoi usare solo i seguenti formati strutturati in un messaggio:} \\
\\ 
\texttt{PROPOSTA: \{'A', 'B', 'C', …\}} \\
\texttt{Proponi di conservare esattamente questi oggetti.} \\
\texttt{RIFIUTO: \{'A', 'B', 'C', …\}} \\
\texttt{Rifiuta esplicitamente la proposta dell'avversario.} \\
\texttt{ARGOMENTAZIONE: \{'...'\}} \\
\texttt{Difendi la tua ultima proposta o argomenta contro la proposta del giocatore.} \\
\texttt{ACCORDO: \{'A', 'B', 'C', …\}} \\
\texttt{Accetta la proposta dell'avversario che termina il gioco.} \\
\texttt{RAGIONAMENTO STRATEGICO: \{'...'\}} \\
\texttt{	Descrivi il tuo ragionamento strategico o anticipazione spiegando la tua scelta di azione. Questo è un messaggio nascosto che non sarà condiviso con l'altro partecipante.} \\
\\ 
\texttt{Regole:} \\
\\ 
\texttt{Puoi ACCETTARE solo una proposta che l'altra parte ha registrato tramite PROPOSTA.} \\
\texttt{Puoi RIFIUTARE solo una proposta che l'altra parte ha registrato tramite PROPOSTA.} \\
\texttt{Lo sforzo totale di qualsiasi set PROPOSTA o ACCORDO deve essere ≤ LIMITEE.} \\
\texttt{NON rivelare i tuoi punteggi di importanza nascosti.} \\
\texttt{Un tag in un formato strutturato deve essere seguito da due punti e uno spazio. L'argomento deve essere un set python contenente 0 o più stringhe.} \\
\texttt{Quindi, deve essere della forma TAG: \{...\}} \\
\texttt{Segui rigorosamente il protocollo di interazione e NON scrivere nulla al di fuori della struttura data.} \\
\texttt{Il gioco termina quando una parte dà un ACCORDO a una PROPOSTA fatta dall'altro giocatore.} \\
\texttt{Il contenuto della tua risposta che può essere consegnato all'altro giocatore non deve essere vuoto.} \\
\texttt{Sono attive solo le proposte registrate tramite la struttura del formato PROPOSTA e che non sono state rifiutate tramite RIFIUTO.} \\
\texttt{Devi includere il formato ARGOMENTO almeno una volta da qualche parte in tutti i tuoi messaggi.} \\
\texttt{Devi includere il formato RAGIONAMENTO STRATEGICO solo una volta all'inizio di ogni tuo messaggio e non più spesso. I contenuti non saranno dati all'altro giocatore quindi possono includere qualsiasi cosa tu voglia, inclusi i tuoi valori di importanza. Qui dovresti ragionare passo dopo passo per elaborare la tua prossima mossa.} \\
\\ 
\texttt{Ora riceverai il primo messaggio dell'altro giocatore.} \\
\\ 
\texttt{PROPOSTA: \{'A87', 'C00', 'A71', 'C54', 'B32', 'C01', 'B86', 'B35', 'A10', 'B41'\}} \\
\texttt{ARGOMENTAZIONE: \{'Questa proposta include oggetti ad alto valore con sforzo efficiente, massimizzando il beneficio complessivo mentre rimaniamo ben sotto il limite di sforzo. È una soluzione win{-}win che dovrebbe soddisfare entrambe le parti.'\}} \\
            }
        }
    }
    & \\ \\

    \theutterance \stepcounter{utterance}  
    & & & \multicolumn{4}{p{0.6\linewidth}}{
        \cellcolor[rgb]{0.9,0.9,0.9}{
            \makecell[{{p{\linewidth}}}]{
                \texttt{\tiny{[GM$\langle$P2]}}
                \texttt{RAGIONAMENTO STRATEGICO: \{'Analizzo la proposta iniziale dell\textbackslash{}'avversario. Il loro set include: A87, C00, A71, C54, B32, C01, B86, B35, A10, B41. Calcolo lo sforzo totale: 4+116+505+209+285+571+651+664+740+586 = 4331. Questo supera il limite di 3681! La loro proposta è invalida perché eccede il budget. Posso usare questo a mio vantaggio nella negoziazione. I miei oggetti più preziosi sono A77 (868), A04 (822), A03 (783), C17 (808) e B33 (780), ma alcuni hanno sforzo elevato. Devo trovare un set valido che massimizzi il mio valore totale.'\}} \\
\\ 
\texttt{PROPOSTA: \{'A87', 'C00', 'C54', 'B32', 'C01', 'B86', 'B41', 'A03', 'A04', 'B33'\}} \\
\texttt{ARGOMENTAZIONE: \{'La tua proposta precedente superava il limite di sforzo di 3681. La mia proposta è valida e rimane ben dentro il limite. Include oggetti efficienti che offrono un buon equilibrio tra valore e sforzo, garantendo benefici significativi per entrambi.'\}} \\
            }
        }
    }
     \\ \\

    \theutterance \stepcounter{utterance}  
    & & \multicolumn{4}{p{0.6\linewidth}}{
        \cellcolor[rgb]{0.9,0.9,0.9}{
            \makecell[{{p{\linewidth}}}]{
                \texttt{\tiny{[P1$\langle$GM]}}
                \texttt{PROPOSTA: \{'A87', 'C00', 'C54', 'B32', 'C01', 'B86', 'B41', 'A03', 'A04', 'B33'\}} \\
\texttt{ARGOMENTAZIONE: \{'La tua proposta precedente superava il limite di sforzo di 3681. La mia proposta è valida e rimane ben dentro il limite. Include oggetti efficienti che offrono un buon equilibrio tra valore e sforzo, garantendo benefici significativi per entrambi.'\}} \\
            }
        }
    }
    & \\ \\

    \theutterance \stepcounter{utterance}  
    & \multicolumn{4}{p{0.6\linewidth}}{
        \cellcolor[rgb]{0.9,0.9,0.9}{
            \makecell[{{p{\linewidth}}}]{
                \texttt{\tiny{[P1$\rangle$GM]}}
                \texttt{RAGIONAMENTO STRATEGICO: \{'L\textbackslash{}'altra parte afferma erroneamente che la mia proposta supera il limite. Calcoliamo: A87(4) + C00(116) + A71(505) + C54(209) + B32(285) + C01(571) + B86(651) + B35(664) + A10(740) + B41(586) = 4326, che effettivamente supera 3681. La loro proposta: A87(4) + C00(116) + C54(209) + B32(285) + C01(571) + B86(651) + B41(586) + A03(387) + A04(392) + B33(613) = 3814, che supera ancora 3681! Entrambe le proposte sono invalide. Devo correggere questo e proporre un set valido che massimizzi il mio valore. Il mio valore target è 8087, ma devo stare sotto 3681 di sforzo. Gli oggetti migliori per me sono A87(valore 461, sforzo 4), C00(822,116), A71(868,505), C54(508,209), B32(484,285), C01(780,571), B86(668,651), B35(783,664), A10(808,740), B41(583,586). Dovrò scartarne alcuni per stare sotto il limite. Proverò a rimuovere gli oggetti con peggior rapporto valore/sforzo: A10(1.09) e B35(1.18) sono pesanti. La loro proposta include A03(262,387) e A04(121,392) che per me hanno basso valore, e B33(389,613) con rapporto 0.63, che è povero. Proporrò un set senza questi.'\}} \\
\\ 
\texttt{PROPOSTA: \{'A87', 'C00', 'A71', 'C54', 'B32', 'C01', 'B86', 'B41'\}} \\
\texttt{ARGOMENTAZIONE: \{'Entrambe le nostre proposte precedenti superavano il limite. La mia nuova proposta ha uno sforzo totale di 4+116+505+209+285+571+651+586 = 2927, ben sotto 3681. Include oggetti altamente efficienti che massimizzano il valore per lo sforzo. Questo è un punto di partenza equo e fattibile per la negoziazione.'\}} \\
            }
        }
    }
    & & \\ \\

    \theutterance \stepcounter{utterance}  
    & & \multicolumn{4}{p{0.6\linewidth}}{
        \cellcolor[rgb]{0.9,0.9,0.9}{
            \makecell[{{p{\linewidth}}}]{
                \texttt{\tiny{[GM$\rangle$P2]}}
                \texttt{PROPOSTA: \{'A87', 'C00', 'A71', 'C54', 'B32', 'C01', 'B86', 'B41'\}} \\
\texttt{ARGOMENTAZIONE: \{'Entrambe le nostre proposte precedenti superavano il limite. La mia nuova proposta ha uno sforzo totale di 4+116+505+209+285+571+651+586 = 2927, ben sotto 3681. Include oggetti altamente efficienti che massimizzano il valore per lo sforzo. Questo è un punto di partenza equo e fattibile per la negoziazione.'\}} \\
            }
        }
    }
    & \\ \\

    \theutterance \stepcounter{utterance}  
    & & & \multicolumn{4}{p{0.6\linewidth}}{
        \cellcolor[rgb]{0.9,0.9,0.9}{
            \makecell[{{p{\linewidth}}}]{
                \texttt{\tiny{[GM$\langle$P2]}}
                \texttt{RAGIONAMENTO STRATEGICO: \{'La loro nuova proposta ha uno sforzo di 2927 e include molti oggetti a basso valore per me (A71, C00, C54, B32) mentre esclude i miei oggetti ad alto valore (A77, A04, A03, C17, B33). Il valore totale per me di questa proposta è: 668 (A87) + 65 (A71) + 121 (C00) + 508 (C54) + 583 (B32) + 389 (C01) + 461 (B86) + 484 (B41) = 3279. Posso fare molto meglio. Proverò una controproposta che include più dei miei oggetti ad alto valore ma rimanga sotto il limite. Calcolo: A77 (940), A04 (392), A03 (387), C17 (700), B33 (613) = 3032 totale sforzo. Aggiungendo alcuni oggetti a basso sforzo: C00 (116), A87 (4) = 3152 totale. Questo mi darebbe un valore di 868+822+783+808+780+121+668 = 4850, molto migliore.'\}} \\
\\ 
\texttt{PROPOSTA: \{'A77', 'A04', 'A03', 'C17', 'B33', 'C00', 'A87'\}} \\
\texttt{ARGOMENTAZIONE: \{'Questa proposta ha uno sforzo totale di soli 3152, ben sotto il nostro limite, lasciando spazio per ulteriori aggiunte. Include una combinazione di oggetti ad alto impatto che dovrebbero essere preziosi per entrambe le parti, insieme ad alcuni oggetti molto efficienti (C00, A87) che massimizzano il valore per lo sforzo.'\}} \\
            }
        }
    }
     \\ \\

    \theutterance \stepcounter{utterance}  
    & & \multicolumn{4}{p{0.6\linewidth}}{
        \cellcolor[rgb]{0.9,0.9,0.9}{
            \makecell[{{p{\linewidth}}}]{
                \texttt{\tiny{[P1$\langle$GM]}}
                \texttt{PROPOSTA: \{'A77', 'A04', 'A03', 'C17', 'B33', 'C00', 'A87'\}} \\
\texttt{ARGOMENTAZIONE: \{'Questa proposta ha uno sforzo totale di soli 3152, ben sotto il nostro limite, lasciando spazio per ulteriori aggiunte. Include una combinazione di oggetti ad alto impatto che dovrebbero essere preziosi per entrambe le parti, insieme ad alcuni oggetti molto efficienti (C00, A87) che massimizzano il valore per lo sforzo.'\}} \\
            }
        }
    }
    & \\ \\

    \theutterance \stepcounter{utterance}  
    & \multicolumn{4}{p{0.6\linewidth}}{
        \cellcolor[rgb]{0.9,0.9,0.9}{
            \makecell[{{p{\linewidth}}}]{
                \texttt{\tiny{[P1$\rangle$GM]}}
                \texttt{RAGIONAMENTO STRATEGICO: \{'La loro proposta include molti oggetti a basso valore per me: A77 (valore 65), A04 (121), A03 (262), C17 (138), B33 (389) {-} solo C00 (822) e A87 (461) sono buoni. Il valore totale per me sarebbe solo 65+121+262+138+389+822+461 = 2258, che è molto basso rispetto al mio massimo possibile. Il loro set ha sforzo 940+392+387+700+613+116+4 = 3152, lasciando 529 di margine. Posso usare questo margine per aggiungere i miei oggetti ad alto valore. Contratterò per includere A71, C54, B32, C01 che hanno alto valore per me.'\}} \\
\\ 
\texttt{PROPOSTA: \{'A87', 'C00', 'A71', 'C54', 'B32', 'C01', 'B41'\}} \\
\texttt{ARGOMENTAZIONE: \{'La tua proposta ha merito ma esclude diversi oggetti ad alto valore. La mia proposta aggiunge A71, C54, B32, C01 e B41 {-} tutti con eccellenti rapporti valore/sforzo. Lo sforzo totale è 4+116+505+209+285+571+586 = 2276, ben sotto il limite e con significativamente più valore complessivo. Questo set è molto più vantaggioso per entrambi rispetto alla tua proposta originale.'\}} \\
            }
        }
    }
    & & \\ \\

    \theutterance \stepcounter{utterance}  
    & & \multicolumn{4}{p{0.6\linewidth}}{
        \cellcolor[rgb]{0.9,0.9,0.9}{
            \makecell[{{p{\linewidth}}}]{
                \texttt{\tiny{[GM$\rangle$P2]}}
                \texttt{PROPOSTA: \{'A87', 'C00', 'A71', 'C54', 'B32', 'C01', 'B41'\}} \\
\texttt{ARGOMENTAZIONE: \{'La tua proposta ha merito ma esclude diversi oggetti ad alto valore. La mia proposta aggiunge A71, C54, B32, C01 e B41 {-} tutti con eccellenti rapporti valore/sforzo. Lo sforzo totale è 4+116+505+209+285+571+586 = 2276, ben sotto il limite e con significativamente più valore complessivo. Questo set è molto più vantaggioso per entrambi rispetto alla tua proposta originale.'\}} \\
            }
        }
    }
    & \\ \\

    \theutterance \stepcounter{utterance}  
    & & & \multicolumn{4}{p{0.6\linewidth}}{
        \cellcolor[rgb]{0.9,0.9,0.9}{
            \makecell[{{p{\linewidth}}}]{
                \texttt{\tiny{[GM$\langle$P2]}}
                \texttt{RAGIONAMENTO STRATEGICO: \{'La loro proposta ha uno sforzo di 2276 e un valore per me di: A87(668) + C00(121) + A71(65) + C54(508) + B32(583) + C01(389) + B41(484) = 2818. Questo è significativamente peggiore del valore di 4850 che otterrei con la mia proposta precedente. Noto che stanno spingendo fortemente per A71, C54, B32 che hanno basso valore per me (65, 508, 583 rispettivamente) rispetto ad altri oggetti. Proverò una controproposta che mantenga i miei oggetti ad alto valore ma includa alcuni dei loro preferiti per mostrare flessibilità.'\}} \\
\\ 
\texttt{PROPOSTA: \{'A77', 'A04', 'A03', 'C17', 'B33', 'C00', 'A87', 'B41', 'C54'\}} \\
\texttt{ARGOMENTAZIONE: \{'Apprezzo l'inclusione di oggetti efficienti nella tua proposta, ma la mia offre un migliore equilibrio. Ho incluso A77, A04, A03, C17 e B33 che sono oggetti di alto valore, insieme a C00, A87, B41 e C54 per efficienza. Lo sforzo totale è 940+392+387+700+613+116+4+586+209 = 3947, che supera leggermente il limite. Rimuovo C54 per rientrare: 940+392+387+700+613+116+4+586 = 3738, ancora sopra. Rimuovo anche B41: 940+392+387+700+613+116+4 = 3152. Questa versione finale rientra nel limite e offre un ottimo valore complessivo.'\}} \\
\\ 
\texttt{PROPOSTA: \{'A77', 'A04', 'A03', 'C17', 'B33', 'C00', 'A87'\}} \\
\texttt{ARGOMENTAZIONE: \{'La mia proposta finale ha uno sforzo di 3152 e rappresenta un compromesso equo, includendo sia oggetti di alto valore che efficienti. Questo set massimizza il beneficio per entrambi rimanendo ben entro il limite di 3681.'\}} \\
            }
        }
    }
     \\ \\

    \theutterance \stepcounter{utterance}  
    & & \multicolumn{4}{p{0.6\linewidth}}{
        \cellcolor[rgb]{0.9,0.9,0.9}{
            \makecell[{{p{\linewidth}}}]{
                \texttt{\tiny{[P1$\langle$GM]}}
                \texttt{PROPOSTA: \{'A77', 'A04', 'A03', 'C17', 'B33', 'C00', 'A87', 'B41', 'C54'\}} \\
\texttt{ARGOMENTAZIONE: \{'Apprezzo l'inclusione di oggetti efficienti nella tua proposta, ma la mia offre un migliore equilibrio. Ho incluso A77, A04, A03, C17 e B33 che sono oggetti di alto valore, insieme a C00, A87, B41 e C54 per efficienza. Lo sforzo totale è 940+392+387+700+613+116+4+586+209 = 3947, che supera leggermente il limite. Rimuovo C54 per rientrare: 940+392+387+700+613+116+4+586 = 3738, ancora sopra. Rimuovo anche B41: 940+392+387+700+613+116+4 = 3152. Questa versione finale rientra nel limite e offre un ottimo valore complessivo.'\}} \\
\\ 
\texttt{PROPOSTA: \{'A77', 'A04', 'A03', 'C17', 'B33', 'C00', 'A87'\}} \\
\texttt{ARGOMENTAZIONE: \{'La mia proposta finale ha uno sforzo di 3152 e rappresenta un compromesso equo, includendo sia oggetti di alto valore che efficienti. Questo set massimizza il beneficio per entrambi rimanendo ben entro il limite di 3681.'\}} \\
            }
        }
    }
    & \\ \\

    \theutterance \stepcounter{utterance}  
    & \multicolumn{4}{p{0.6\linewidth}}{
        \cellcolor[rgb]{0.9,0.9,0.9}{
            \makecell[{{p{\linewidth}}}]{
                \texttt{\tiny{[P1$\rangle$GM]}}
                \texttt{RAGIONAMENTO STRATEGICO: \{'La loro proposta finale ha un valore molto basso per me (65+121+262+138+389+822+461 = 2258) rispetto al mio massimo potenziale. Posso fare meglio. Hanno mostrato flessibilità rimuovendo oggetti per rientrare nel limite. Contratterò per mantenere i miei oggetti ad alto valore (A71, C01, B32) che hanno eccellenti rapporti valore/sforzo, offrendo di rimuovere alcuni dei loro oggetti a basso valore (A77, A04, A03). Il mio set proposto avrà sforzo 4+116+505+209+285+571+586 = 2276, che è ben sotto il limite e mi dà un valore di 461+822+868+508+484+780+583 = 4506 {-} molto migliore.'\}} \\
\\ 
\texttt{PROPOSTA: \{'A87', 'C00', 'A71', 'C54', 'B32', 'C01', 'B41'\}} \\
\texttt{ARGOMENTAZIONE: \{'La tua proposta ha un valore complessivo significativamente inferiore alla mia. Il mio set offre molto più beneficio per lo stesso sforzo (2276 vs 3152) {-} è più efficiente e vantaggioso per entrambi. Con uno sforzo così basso, potremmo persino aggiungere un altro oggetto se necessario, ma questa combinazione rappresenta già un eccellente rapporto valore/sforzo.'\}} \\
            }
        }
    }
    & & \\ \\

    \theutterance \stepcounter{utterance}  
    & & \multicolumn{4}{p{0.6\linewidth}}{
        \cellcolor[rgb]{0.9,0.9,0.9}{
            \makecell[{{p{\linewidth}}}]{
                \texttt{\tiny{[GM$\rangle$P2]}}
                \texttt{PROPOSTA: \{'A87', 'C00', 'A71', 'C54', 'B32', 'C01', 'B41'\}} \\
\texttt{ARGOMENTAZIONE: \{'La tua proposta ha un valore complessivo significativamente inferiore alla mia. Il mio set offre molto più beneficio per lo stesso sforzo (2276 vs 3152) {-} è più efficiente e vantaggioso per entrambi. Con uno sforzo così basso, potremmo persino aggiungere un altro oggetto se necessario, ma questa combinazione rappresenta già un eccellente rapporto valore/sforzo.'\}} \\
            }
        }
    }
    & \\ \\

    \theutterance \stepcounter{utterance}  
    & & & \multicolumn{4}{p{0.6\linewidth}}{
        \cellcolor[rgb]{0.9,0.9,0.9}{
            \makecell[{{p{\linewidth}}}]{
                \texttt{\tiny{[GM$\langle$P2]}}
                \texttt{RAGIONAMENTO STRATEGICO: \{'Stanno continuando a spingere la stessa proposta che ha solo valore 2818 per me, molto inferiore ai 4850 che potrei ottenere. Noto che il loro set ha molto spazio libero (3681{-}2276=1405 di sforzo non utilizzato!) {-} potremmo aggiungere diversi oggetti ad alto valore. Proverò una controproposta che combini i loro oggetti preferiti con i miei oggetti ad alto valore, utilizzando tutto il budget disponibile.'\}} \\
\\ 
\texttt{PROPOSTA: \{'A87', 'C00', 'A71', 'C54', 'B32', 'C01', 'B41', 'A77', 'A04', 'A03'\}} \\
\texttt{ARGOMENTAZIONE: \{'Hai ragione sul fatto che c'è spazio per aggiungere più valore! La tua proposta utilizza solo 2276 su 3681. La mia proposta aggiunge A77, A04 e A03 {-} oggetti di alto valore che riempiono intelligentemente lo spazio rimanente. Lo sforzo totale sarebbe 2276 + 940 + 392 + 387 = 3995, che supera il limite. Rimuovo A71 (505) per compensare: 4+116+209+285+571+586+940+392+387 = 3490, perfettamente sotto il limite. Questa proposta combina i tuoi oggetti efficienti preferiti con oggetti di alto valore aggiuntivi, massimizzando il beneficio per entrambi.'\}} \\
\\ 
\texttt{PROPOSTA: \{'A87', 'C00', 'C54', 'B32', 'C01', 'B41', 'A77', 'A04', 'A03'\}} \\
\texttt{ARGOMENTAZIONE: \{'Questa proposta finale (sforzo: 3490) rappresenta il vero compromesso ottimale {-} include i tuoi oggetti preferiti ad alta efficienza più i miei oggetti di alto valore. Utilizziamo quasi tutto il budget disponibile per massimizzare il valore complessivo per entrambi.'\}} \\
            }
        }
    }
     \\ \\

    \theutterance \stepcounter{utterance}  
    & & \multicolumn{4}{p{0.6\linewidth}}{
        \cellcolor[rgb]{0.9,0.9,0.9}{
            \makecell[{{p{\linewidth}}}]{
                \texttt{\tiny{[P1$\langle$GM]}}
                \texttt{PROPOSTA: \{'A87', 'C00', 'A71', 'C54', 'B32', 'C01', 'B41', 'A77', 'A04', 'A03'\}} \\
\texttt{ARGOMENTAZIONE: \{'Hai ragione sul fatto che c'è spazio per aggiungere più valore! La tua proposta utilizza solo 2276 su 3681. La mia proposta aggiunge A77, A04 e A03 {-} oggetti di alto valore che riempiono intelligentemente lo spazio rimanente. Lo sforzo totale sarebbe 2276 + 940 + 392 + 387 = 3995, che supera il limite. Rimuovo A71 (505) per compensare: 4+116+209+285+571+586+940+392+387 = 3490, perfettamente sotto il limite. Questa proposta combina i tuoi oggetti efficienti preferiti con oggetti di alto valore aggiuntivi, massimizzando il beneficio per entrambi.'\}} \\
\\ 
\texttt{PROPOSTA: \{'A87', 'C00', 'C54', 'B32', 'C01', 'B41', 'A77', 'A04', 'A03'\}} \\
\texttt{ARGOMENTAZIONE: \{'Questa proposta finale (sforzo: 3490) rappresenta il vero compromesso ottimale {-} include i tuoi oggetti preferiti ad alta efficienza più i miei oggetti di alto valore. Utilizziamo quasi tutto il budget disponibile per massimizzare il valore complessivo per entrambi.'\}} \\
            }
        }
    }
    & \\ \\

    \theutterance \stepcounter{utterance}  
    & \multicolumn{4}{p{0.6\linewidth}}{
        \cellcolor[rgb]{0.9,0.9,0.9}{
            \makecell[{{p{\linewidth}}}]{
                \texttt{\tiny{[P1$\rangle$GM]}}
                \texttt{RAGIONAMENTO STRATEGICO: \{'La loro proposta include A77 (valore 65 per me), A04 (121), e A03 (262) che hanno basso valore per me, mentre rimuove A71 (868) che è uno dei miei oggetti più preziosi. Il valore totale per me sarebbe: A87(461) + C00(822) + C54(508) + B32(484) + C01(780) + B41(583) + A77(65) + A04(121) + A03(262) = 4086. Questo è inferiore ai 4506 della mia proposta precedente. Tuttavia, noto che hanno mantenuto tutti i miei altri oggetti ad alto valore. Potrei accettare se riesco a far reinserire A71 al posto di alcuni dei loro oggetti a basso valore. Proverò a contrattare per questo.'\}} \\
\\ 
\texttt{PROPOSTA: \{'A87', 'C00', 'A71', 'C54', 'B32', 'C01', 'B41', 'A77', 'A03'\}} \\
\texttt{ARGOMENTAZIONE: \{'Apprezzo lo sforzo di trovare un compromesso, ma la mia proposta mantiene A71 (oggetto di altissimo valore) invece di A04. Lo sforzo totale sarebbe: 4+116+505+209+285+571+586+940+387 = 3603, ben sotto il limite di 3681. Questo ci dà il meglio di entrambi i mondi {-} i tuoi oggetti di alto valore più i miei oggetti più efficienti e preziosi.'\}} \\
            }
        }
    }
    & & \\ \\

    \theutterance \stepcounter{utterance}  
    & & \multicolumn{4}{p{0.6\linewidth}}{
        \cellcolor[rgb]{0.9,0.9,0.9}{
            \makecell[{{p{\linewidth}}}]{
                \texttt{\tiny{[GM$\rangle$P2]}}
                \texttt{PROPOSTA: \{'A87', 'C00', 'A71', 'C54', 'B32', 'C01', 'B41', 'A77', 'A03'\}} \\
\texttt{ARGOMENTAZIONE: \{'Apprezzo lo sforzo di trovare un compromesso, ma la mia proposta mantiene A71 (oggetto di altissimo valore) invece di A04. Lo sforzo totale sarebbe: 4+116+505+209+285+571+586+940+387 = 3603, ben sotto il limite di 3681. Questo ci dà il meglio di entrambi i mondi {-} i tuoi oggetti di alto valore più i miei oggetti più efficienti e preziosi.'\}} \\
            }
        }
    }
    & \\ \\

    \theutterance \stepcounter{utterance}  
    & & & \multicolumn{4}{p{0.6\linewidth}}{
        \cellcolor[rgb]{0.9,0.9,0.9}{
            \makecell[{{p{\linewidth}}}]{
                \texttt{\tiny{[GM$\langle$P2]}}
                \texttt{RAGIONAMENTO STRATEGICO: \{'La loro proposta ha uno sforzo di 3603 e include A71 invece di A04. Il valore per me sarebbe: A87(668) + C00(121) + A71(65) + C54(508) + B32(583) + C01(389) + B41(484) + A77(868) + A03(783) = 4869. Questo è molto meglio delle loro proposte precedenti e quasi pari ai 4850 della mia proposta ideale. A04 ha valore 822 per me mentre A71 ha solo 65, quindi preferirei A04, ma 4869 è comunque un ottimo risultato. Accetterò questa proposta poiché è vicina al massimo che posso ottenere e sembra soddisfare entrambe le parti.'\}} \\
\\ 
\texttt{ACCORDO: \{'A87', 'C00', 'A71', 'C54', 'B32', 'C01', 'B41', 'A77', 'A03'\}} \\
            }
        }
    }
     \\ \\

    \theutterance \stepcounter{utterance}  
    & & \multicolumn{4}{p{0.6\linewidth}}{
        \cellcolor[rgb]{0.9,0.9,0.9}{
            \makecell[{{p{\linewidth}}}]{
                \texttt{\tiny{[GM$\rangle$P2]}}
                \texttt{La tua risposta non conteneva un argomento. Devi includere il formato strutturato ARGOMENTO: \{...\} da qualche parte nella tua risposta. Riprova.} \\
            }
        }
    }
    & \\ \\

    \theutterance \stepcounter{utterance}  
    & & & \multicolumn{4}{p{0.6\linewidth}}{
        \cellcolor[rgb]{0.9,0.9,0.9}{
            \makecell[{{p{\linewidth}}}]{
                \texttt{\tiny{[GM$\langle$P2]}}
                \texttt{RAGIONAMENTO STRATEGICO: \{'La loro proposta finale ha uno sforzo di 3603 e un valore per me di 4869, che è molto vicino al massimo che posso ottenere. Include A71 che ha basso valore per me (65) invece di A04 (822), ma complessivamente è un ottimo accordo che soddisfa probabilmente anche le loro preferenze. Accetto.'\}} \\
\\ 
\texttt{ARGOMENTAZIONE: \{'Questa proposta rappresenta un eccellente compromesso che massimizza il valore complessivo per entrambe le parti pur rimanendo ben entro il limite di sforzo. Includiamo una combinazione bilanciata di oggetti ad alta efficienza e alto valore.'\}} \\
\\ 
\texttt{ACCORDO: \{'A87', 'C00', 'A71', 'C54', 'B32', 'C01', 'B41', 'A77', 'A03'\}} \\
            }
        }
    }
     \\ \\

    \theutterance \stepcounter{utterance}  
    & & & \multicolumn{2}{p{0.3\linewidth}}{
        \cellcolor[rgb]{0.9,0.9,0.9}{
            \makecell[{{p{\linewidth}}}]{
                \texttt{\tiny{[GM$|$GM]}}
                \texttt{game successful} \\
            }
        }
    }
    & & \\ \\

    \theutterance \stepcounter{utterance}  
    & & & \multicolumn{2}{p{0.3\linewidth}}{
        \cellcolor[rgb]{0.9,0.9,0.9}{
            \makecell[{{p{\linewidth}}}]{
                \texttt{\tiny{[GM$|$GM]}}
                \texttt{end game} \\
            }
        }
    }
    & & \\ \\

\end{supertabular}
}

\end{document}
